\chapter{uafR Table and Script Guide}

\section{High-Value categorate Tables}
\begin{longtable}{p{0.34\textwidth}p{0.54\textwidth}}
\toprule
\textbf{Table} & \textbf{Research use}\\
\midrule
\endfirsthead
\toprule
\textbf{Table} & \textbf{Research use}\\
\midrule
\endhead
\texttt{PubChemIdentity} & Check CID matches, match status, and source URLs.\\
\texttt{PubChemProperties} & Compare formula, exact mass, molecular weight, and descriptor fields when available.\\
\texttt{SourceCoverage} & Identify which databases contributed evidence for each chemical.\\
\texttt{ChemicalTerms} & Inspect long-form extracted terms before committing to variables.\\
\texttt{ChemicalTraits} & Use cross-source discrete traits with confidence and evidence fields.\\
\texttt{ChemicalTraitOntology} & Use whitelisted ontology-style traits grouped by research domain.\\
\texttt{ChemicalTraitMatrix} & Analyze chemicals by binary, count, or confidence-weighted trait columns.\\
\texttt{ChemicalTraitEvidence} & Trace matrix columns back to source evidence.\\
\texttt{ChemicalTraitReport} & Review one compact row per compound for high-level interpretation.\\
\texttt{ChemicalMeasurementSummary} & Summarize extracted measured or computed properties.\\
\texttt{ChemicalHazards} & Inspect source-backed safety and hazard values.\\
\texttt{ChemicalOccurrences} & Review organism or natural product occurrence evidence.\\
\texttt{ChemicalPathwayRoles} & Interpret KEGG pathway context and biochemical roles.\\
\texttt{KEGGReactionParticipants} & Examine reaction participation and related biochemical context.\\
\texttt{DataDictionary} & Confirm expected table schemas and field meaning.\\
\texttt{TableQuality} & Audit missing columns, duplicate keys, and completeness.\\
\texttt{ValidationSummary} & Start every review here before interpreting output.\\
\bottomrule
\end{longtable}

\section{Training Scripts}
\begin{tabularx}{\textwidth}{p{0.34\textwidth}Y}
\toprule
\textbf{Script} & \textbf{Purpose}\\
\midrule
\path{00_install_check.R} & Confirms R, package, dependencies, and example data.\\
\path{01_project_setup.R} & Creates a reproducible student project folder.\\
\path{04_core_workflow.R} & Runs the core GC-MS example workflow.\\
\path{05_categorate_research.R} & Runs a small research enrichment smoke test.\\
\path{06_trait_matrix_analysis.R} & Builds selected trait matrices and evidence tables.\\
\path{07_visualization.R} & Creates example coverage and trait figures.\\
\path{09_reproducibility_check.R} & Audits a student project folder before final submission.\\
\path{run_training_smoke_tests.R} & Runs lightweight checks across the training scripts.\\
\bottomrule
\end{tabularx}

\section{Student Bundle Scripts}
\begin{tabularx}{\textwidth}{p{0.38\textwidth}Y}
\toprule
\textbf{Script} & \textbf{Purpose}\\
\midrule
\path{tools/build_student_bundle.R} & Builds the private student bundle folder and zip archive from the repository root.\\
\path{START_HERE.md} & Gives the first student checklist for setup, installation, acceptance testing, and opening the manual.\\
\path{preflight_check.R} & Checks R, RStudio, internet access, bundle contents, library write access, and required dependencies before installation.\\
\path{install_uafR_from_bundle.R} & Installs dependencies and the local uafR source archive from the unzipped student bundle.\\
\path{update_uafR_from_bundle.R} & Reinstalls uafR from a newer bundle and verifies the result.\\
\path{verify_uafR_install.R} & Confirms package loading, dependencies, bundled datasets, and saved standard categorate data without live web queries.\\
\path{run_student_acceptance_test.R} & Runs package verification, the offline training workflow, and an optional live database smoke test.\\
\path{examples/test_install.R} & Reruns the acceptance test from the bundle examples folder.\\
\bottomrule
\end{tabularx}

\section{Choosing Detail Levels}
\begin{itemize}
  \item Use \texttt{detail = "standard"} for the legacy eight-table output.
  \item Use \texttt{detail = "research"} for most student projects. It gives rich tables while keeping the query smaller than full assay detail.
  \item Use \texttt{detail = "full"} only when the project needs broad assay details or deeper PubChem annotations.
\end{itemize}

\section{Recommended Smoke Test Settings}
\begin{rcode}
categorate(
  compounds = c("aspirin", "caffeine"),
  chemical_library = library_data,
  input_format = "wide",
  detail = "research",
  cache = TRUE,
  throttle = 0.2,
  assay_detail_limit = 0,
  trait_matrix_profile = "core",
  trait_matrix_min_confidence = "medium",
  trait_matrix_max_traits = 100
)
\end{rcode}
